\begin{summarybox}
\paragraph{核心：在同一区间内，用导数控制两点间的高度差}
在连续、内部可导的前提下：
\[
f'>0\Rightarrow\text{严格递增},\qquad f'\ge0\Longleftrightarrow\text{不减}.
\]
\[
\bigl[f'\ge0\ \text{且无常值子区间}\bigr]\Longrightarrow\text{严格递增}.
\]
递减方向将符号反向。最后检查：定义域是否分段？端点是否连续？是否误把充分条件当必要条件？零导数点是否真的形成常值区间？
\end{summarybox}
